%==============================================================================
% CONFIG: LAYOUT, TIPOGRAFI, DAN SPASI
%==============================================================================

% Font: Times New Roman (fallback ke TeX Gyre Termes bila font sistem belum terpasang)
\IfFontExistsTF{Times New Roman}{%
  \setmainfont{Times New Roman}%
}{%
  \setmainfont{TeX Gyre Termes}%
}

% Tipografi: Times New Roman 12 pt, rata kiri-kanan, spasi 1,5 (setara Microsoft Word)
%   Word "1.5 lines" TNR 12 pt = 20,7 pt  ->  skala 1,43 dari baselineskip LaTeX
%   Word "single"     TNR 12 pt = 13,8 pt  ->  skala 0,952
\newcommand{\spasisatu}{\setstretch{0.952}}
\newcommand{\spasisetengah}{\setstretch{1.43}}
\AtBeginDocument{\spasisetengah}

% Indentasi & spasi antar paragraf
\setlength{\parindent}{1.27cm}   % First-line indent 720 twips (1,27 cm)
\setlength{\parskip}{10pt}       % Spacing after paragraf 200 twips (10 pt)
\setlength{\emergencystretch}{2em}

% Label penamaan bahasa Indonesia
\renewcommand{\figurename}{Gambar}
\renewcommand{\tablename}{Tabel}
\renewcommand{\contentsname}{Daftar Isi}
